% =====================================================================
% Bagian 7 — Kesimpulan dan Pengembangan Lanjutan  (target 0,5 halaman)
% =====================================================================

\section{Kesimpulan dan Pengembangan Lanjutan}

Penelitian ini membangun sistem deteksi dini karhutla gambut berjenjang dengan
inti model fusi RGB--termal berbasis \emph{cross-attention} yang dilatih dengan
\emph{modality dropout}, dan mengujinya melalui protokol evaluasi ketahanan
berbasis kurva degradasi kontinu.

\textbf{Yang berhasil ditunjukkan.} Fusi lebih tahan terhadap degradasi dibanding
jalur RGB-saja secara konsisten di ketiga partisi data, dan keunggulannya atas
jalur termal-saja signifikan pada partisi latih dan validasi ($p = 0{,}00098$ dan
$p = 0{,}0078$, McNemar berpasangan). Tidak ada satu pun kejadian api yang
terlewat pada seluruh level degradasi. Mekanisme \emph{gating} berjalan dan
skornya sampai ke operator. Model menunjukkan generalisasi \emph{zero-shot} 92\%
ke dataset lain dan berjalan pada 117,5~ms per bingkai di CPU laptop. Seluruh
komponen terintegrasi menjadi sistem terpasang.

\textbf{Yang tidak berhasil ditunjukkan}, dilaporkan dengan bobot yang sama.
Fusi \emph{tidak} mengungguli jalur RGB-saja pada kondisi bersih, keduanya
menghasilkan prediksi identik pada setiap sampel. \emph{Gating} tidak
meningkatkan akurasi. Kelas \texttt{fire\_no\_smoke} praktis absen dari data
latih dan tidak pernah diprediksi, sehingga sistem tidak dapat mengenali api
tanpa asap, justru rezim yang paling dekat dengan pembakaran membara.
Kuantisasi INT8 menurunkan ukuran model 68,8\% tetapi memperlambat inferensi.
Peringkat ketahanan berbalik pada partisi uji, dan penyebab paling mungkin
adalah operator degradasi termal yang tidak cukup kuat, bukan ketahanan sejati.

Kami menilai temuan negatif ini bukan kelemahan laporan melainkan hasilnya.
Bahwa evaluasi kondisi bersih pada tolok ukur ini tidak memiliki daya pisah sama
sekali adalah informasi yang berguna bagi siapa pun yang hendak membandingkan
arsitektur fusi pada data serupa.

\textbf{Pengembangan lanjutan}, diurutkan menurut kepentingan bagi validitas
klaim: pengumpulan data gambut tropis berlabel, selama tidak ada satu pun
bingkai gambut Indonesia, seluruh klaim tetap berstatus proksi, dan himpunan uji
sekecil apa pun bernilai lebih besar daripada perbaikan arsitektur mana pun;
kalibrasi operator degradasi relatif terhadap margin keputusan model, agar kurva
ketahanan mengukur model dan bukan kelemahan operatornya; pengulangan
multi-benih dan baseline eksternal, yang menutup dua ancaman validitas terbesar
yang tersisa; data teranotasi untuk rezim api-tanpa-asap agar kelas yang kini
mati dapat dipelajari; serta masukan variabel lingkungan sesuai keterbatasan
yang diidentifikasi narasumber ahli dan pengujian pada perangkat tepi fisik
dengan penyebutan model perangkat yang spesifik.
